%% =====================================================================
%%  B1-T-02-02 -- what B1 contributes to "Indifference to Being Liked"
%%  -------------------------------------------------------------------
%%  LAYER A: APPEND-ONLY. Written once, then left alone. Material found
%%  only here sits in the `unique` slot and is protected by rule R10.
%% =====================================================================
%% @kind: source
%% @id: B1-T-02-02
%% @book: B1
%% @topic: T-02-02
%% @locator: Tactic 5, pp. 59--70; Tactic 2, pp. 19--30
%% @status: distilled
%% @unique: yes
%% @integrated: yes
%% @updated: 2026-09-19

\dossier{B1}{T-02-02}{\bookshort{B1}}{Tactic 5, pp. 59--70; Tactic 2, pp. 19--30}

\begin{distilled}
  \item Retaliation is argued against on effectiveness grounds: harsh words produce more harsh words, and an ego injury creates a resentful enemy who is expensive later.
  \item The clever putdown is identified as a fantasy imported from film; in practice it raises resistance rather than clearing it.
  \item The meekness tactic is paired with concealing how much you know --- appearing less formidable makes people talk.
\end{distilled}

\begin{unique}
  \item The explicit anti-retaliation argument, including the observation that winning a verbal exchange produces a worse net outcome than losing it, because the loser becomes a persistent adversary.
\end{unique}
